\subsection{Hilbert 空间的情形}

在本节中，我们考虑基域为 K = R 或 C 的 Hilbert 空间。我们用 \(\langle x_{1}|x_{2}\rangle\) 表示 Hilbert 空间 E 中两个向量 \(x_{1}\) 与 \(x_{2}\) 的内积。

若 E 是复 Hilbert 空间，则 Banach 代数 \(\mathcal{L}(\mathrm{E})\) 赋以对合 \(u\mapsto u^{*}\) 后是一个星代数（I，第 102 页，例 1）。特别地，若 \(u\in \mathcal{L}(\mathrm{E})\) ，则有 \(\varrho (u^{*}) = \varrho (u)\) 及 \(\mathrm{Sp}(u^{*}) = \overline{\mathrm{Sp}(u)}\) 。

设 \(u \in \mathcal{L}(\mathrm{E})\) 为正规自同态，\(\lambda \in C\) 。u 的对应于 \(\lambda\) 的特征子空间与伴随 \(u^{*}\) 的对应于 \(\overline{\lambda}\) 的特征子空间重合（EVT，V，第 43 页，命题 8 的推论）。然而，当 u 不是正规自同态时，这些空间一般并不重合（I，第 188 页，习题 3）。

设 \(\lambda\) 与 \(\mu\) 为复数，满足 \(\lambda \neq \mu\) 。设 x 为 u 的对应于 \(\lambda\) 的特征向量，y 为 u 的对应于 \(\mu\) 的特征向量。则 \(u^{*}(x) = \overline{\lambda}x\) ，从而

\[
\mu \langle x | y \rangle = \langle x | u (y) \rangle = \langle u ^ {*} (x) | y \rangle = \langle \overline {{\lambda}} x | y \rangle = \lambda \langle x | y \rangle .
\]

因此 \(\langle x|y\rangle = 0\) ：u 的诸特征子空间两两正交。此外，对每个 \(\lambda \in C\) ，u 的对应于 \(\lambda\) 的特征子空间与 u 的对应于 \(\lambda\) 的准素子空间重合，后者即（LIE，VII，§1，第 1 节）对 \(k \in N\) 取的诸 \((u - \lambda1_{\mathrm{E}})^{k}\) 的核的并（EVT，V，第 43 页，命题 8 的推论）。

引理 2.　设 E 为复 Hilbert 空间，u 为 E 的正规自同态。设 \((\mathrm{E}_{i})_{i\in\mathrm{I}}\) 为 E 的在 u 下稳定且两两正交的闭子空间的有限族，满足 \(E=\bigoplus_{i\in I}E_{i}\) 。对每个 \(i\in I\) ，用 \(u_{i}\) 表示 u 在 \(E_{i}\) 上诱导的自同态。我们有 \(\mathrm{Sp}(u)=\bigcup_{i\in I}\mathrm{Sp}(u_{i})\) ，且对每个 \(f\in\mathcal{C}(\mathrm{Sp}(u))\) ，自同态 \(f(u)\) 使诸空间 \(E_{i}\) 保持稳定，且 \(f(u)\) 与 \(f(u_{i})\) 在 \(E_{i}\) 上一致。

其证明仿照 I，第 128 页，引理 1 的证明，并用到 I，第 110 页，注记 6 与 I，第 112 页，命题 8。

命题 4.　设 E 为复 Hilbert 空间，u 为 E 的正规自同态。对每个函数 \(f \in \mathcal{C}(\mathrm{Sp}(u))\) 及每个 \(\lambda \in \mathbf{C}\) ，有

\[
\operatorname{Ker} (u - \lambda 1 _ {\mathrm{E}}) \subset \operatorname{Ker} (f (u) - f (\lambda) 1 _ {\mathrm{E}}).
\]

其证明与 I，第 128 页，命题 1 的证明类似；我们复述其论证。同前所引入的代数 A 在此是 \(\mathcal{L}(\mathrm{E})\) 的一个幺星子代数（EVT，V，第 43 页，推论）。因此它包含由 \(u\) 生成的幺星子代数 B，后者是交换的。映射 \(\chi: \mathrm{B} \to \mathbf{C}\) 把 \(v\) 对应到 \(v\) 相对于 \(x\) 的特征值，它是 B 的一个特征标，且 \(\chi(u) = \lambda\) 。对每个 \(f \in \mathcal{C}(\mathrm{Sp}(u))\) ，由 I，第 112 页，命题 8，有 \(f(u) \in \mathrm{B}\) 且 \(\chi(f(u)) = f(\chi(u)) = f(\lambda)\) ，断言得证。

引理 3.　设 E 为 Hilbert 空间，\(p \in \mathcal{L}(\mathrm{E})\) 为一个投影。下列断言等价：

\begin{enumerate}
\def\labelenumi{(\roman{enumi})}
\item
  投影 p 是正交投影，即 \(\operatorname{Ker}(p)=\operatorname{Im}(p)^{\circ}\) （EVT，V，第 13 页）；
\item
  投影 p 是 Hermite 的；
\item
  投影 p 是正规的；
\item
  我们有 \(\operatorname{Ker}(p) \subset \operatorname{Ker}(p^*)\) ；
\item
  我们有 \(\operatorname{Im}(p) \subset \operatorname{Im}(p^*)\) ；
\item
  投影 p 是正的；
\item
  我们有 \(\| p\| \leqslant 1\) 。
\end{enumerate}

首先回忆 \(\operatorname{Ker}(p^{*}) = \operatorname{Im}(p)^{\circ}\) 与 \(\operatorname{Ker}(p) = \operatorname{Im}(p^{*})^{\circ}\) （EVT，V，第 41 页，命题 4）。此外，p（相应地，\(p^{*}\) ）的像是闭的，因为它与 1 − p（相应地，\(1 - p^{*}\) ）的核重合。于是我们有

\[
\operatorname{Im} (p) = \operatorname{Ker} (p ^ {*}) ^ {\circ}, \qquad \operatorname{Im} (p ^ {*}) = \operatorname{Ker} (p) ^ {\circ}.\tag{1}
\]

\begin{enumerate}
\def\labelenumi{(\roman{enumi})}
\item
  \(\Longrightarrow\) (ii)：\(p^{*}\) 是一个投影，其核为 \(\operatorname{Im}(p)^{\circ} = \operatorname{Ker}(p)\) ，其像为 \(\operatorname{Im}(p^{*})^{\circ} = \operatorname{Ker}(p) = \operatorname{Im}(p)^{\circ}\) ；因此 \(p^{*} = p\) 。
\item
  \(\Longrightarrow\) (iii)：因为每个 Hermite 自同态都是正规的。
\item
  \(\Longrightarrow\) (iv)：由于 \(p\) 是正规的，我们有 \(\| p(x)\|^2 = \| p^*(x)\|^2\) 对 E 中所有 \(x\) 成立（EVT，V，第 43 页，命题 7），由此得到所要的包含关系。
\item
  \(\Longrightarrow\) (v) 由上文的等式组 (1) 得出。
\item
  \(\Longrightarrow\) (vi)：对所有 \(x \in E\) ，有 \(p(x) \in \operatorname{Im}(p^{*})\) ，从而 \(\langle p(x)|x\rangle = \langle p^{*}(p(x))|x\rangle = \|p(x)\|^{2} \geqslant 0\) 。
\item
  \(\Longrightarrow\) (vii)：设 \(x \in E\) ，\(y = x - p(x) \in \operatorname{Ker}(p)\) 。对每个 \(t \in R\) ，由假设
\end{enumerate}

\[
\langle x + t y | p (x) \rangle = \langle x + t y | p (x + t y) \rangle \geqslant 0,
\]

这仅当 \(\langle y|p(x)\rangle = 0\) 时才可能。于是

\[
\| p (x) \| ^ {2} = \langle x | p (x) \rangle \leqslant \| x \| \| p (x) \|,
\]

因而 \(\|p\|\leqslant1\) 。

\begin{enumerate}
\def\labelenumi{(\roman{enumi})}
\setcounter{enumi}{6}
\tightlist
\item
  \(\Longrightarrow\) (i)：设 \(y \in \operatorname{Im}(p)\) ；设 \(z\) 为 \(y\) 在 \(\operatorname{Ker}(p)^{\circ}\) 上的正交投影，并记 \(x = y - z \in \operatorname{Ker}(p)\) 。我们有 \(p(z) = p(y) = y\) ，故由假设得 \(\|y\| \leqslant \|z\|\) 。但由于 \(x\) 与 \(z\) 正交，我们有 \(\|y\|^2 = \|x\|^2 + \|z\|^2\) ，由此得 \(\|x\| = 0\) ，即 \(y = z\) 。于是 \(\operatorname{Im}(p) \subset \operatorname{Ker}(p)^{\circ}\) 。又由于 \(\|p^*\| = \|p\| \leqslant 1\) ，同样有 \(\operatorname{Im}(p^*) \subset \operatorname{Ker}(p^*)^{\circ}\) ，再由 (1) 便得到反向包含关系。
\end{enumerate}

命题 5.　设 E 为复 Hilbert 空间，u 为 E 的正规自同态。

\begin{enumerate}
\def\labelenumi{\alph{enumi})}
\item
  对 u 的谱的每个既开又闭子集 H，谱投影 \(e_{\mathrm{H}}(u)\) 是一个正交投影，其核是谱投影 \(e_{\mathrm{Sp}(u)-\mathrm{H}}(u)\) 的像；
\item
  若 \(H_{1}\) 与 \(H_{2}\) 是 u 的谱的不相交的既开又闭子集，则谱子空间 \(E_{H_{1}}\) 与 \(E_{H_{2}}\) 正交；
\item
  若 \(\lambda \in \mathbf{C}\) 是 \(u\) 的谱的孤立点，则 \(\lambda\) 是 \(u\) 的特征值，且谱投影 \(e_{\lambda}(u)\) 的像是 \(u\) 的对应于 \(\lambda\) 的特征子空间。
\end{enumerate}

我们证明 \(a\) ）。由于全纯函数演算与连续函数演算相容（I，第 111 页，推论 1），我们有 \(e_{\mathrm{H}}(u) = \varphi_{\mathrm{H}}(u)\) ，其中 \(\varphi_{\mathrm{H}} \in \mathcal{C}(\mathrm{Sp}(u))\) 是 H 的特征函数。由此推出 \(e_{\mathrm{H}}(u)^{*} = \overline{\varphi}_{\mathrm{H}}(u) = \varphi_{\mathrm{H}}(u) = e_{\mathrm{H}}(u)\) ，故 \(e_{\mathrm{H}}(u)\) 是正交投影（引理 3，(ii)）。它的核是投影 \(1 - e_{\mathrm{H}}(u) = e_{\mathrm{Sp}(u) - \mathrm{H}}(u)\) 的像。

我们证明 \(b\) ）。特征函数 \(\varphi_{\mathrm{H}_{1}}\) 与 \(\varphi_{\mathrm{H}_{2}}\) （\(\mathrm{H}_{1}\) 与 \(\mathrm{H}_{2}\) 在 \(\mathrm{Sp}(u)\) 上的特征函数）都是连续的，且其乘积为零，由此推出 \(e_{\mathrm{H}_{1}}(u)\circ e_{\mathrm{H}_{2}}(u)=e_{\mathrm{H}_{2}}(u)\circ e_{\mathrm{H}_{1}}(u)=0\) 。由它们便得到包含关系 \(\mathrm{E}_{\mathrm{H}_{2}}(u)\subset\mathrm{E}_{\mathrm{H}_{1}}(u)^{\circ}\) 与 \(\mathrm{E}_{\mathrm{H}_{1}}(u)\subset\mathrm{E}_{\mathrm{H}_{2}}(u)^{\circ}\) 。

最后证明断言 \(c\) ）。特征函数 \(\varphi_{\lambda}\) （\(\{\lambda\}\) 的特征函数）在 \(\mathrm{Sp}(u)\) 上连续且非零；它满足 \((z - \lambda)\varphi_{\lambda}(z) = 0\) 对所有 \(z \in \mathrm{Sp}(u)\) 成立。因此我们有 \((u - \lambda 1_{\mathrm{E}})\varphi_{\lambda}(u) = 0\) 。于是 \(\varphi_{\lambda}(u)\) 的非零像包含在 \(u\) 的对应于 \(\lambda\) 的特征子空间中。由于我们有 \(e_{\lambda}(u) = \varphi_{\lambda}(u)\) ，且 \(e_{\lambda}(u)\) 的像包含 \(u\) 的对应于 \(\lambda\) 的特征子空间，断言由此得证。

引理 4.　设 E 为 Hilbert 空间，u 为 E 的正规自同态。设 F 为 E 的一个闭子空间，它包含 u 的特征向量的一个全子集。则 \(\mathrm{F}^{\circ}\) 在 u 下稳定，且自同态 \(\widetilde{u}\) （即由 u 诱导的 \(\mathrm{F}^{\circ}\) 的自同态）是正规的。

由于 \(u\) 是正规的，\(u\) 的每个特征向量也是 \(u^{*}\) 的特征向量（EVT，V，第 43 页，推论）。因此假设蕴涵 F 在 \(u\) 与 \(u^{*}\) 下稳定。于是由 EVT，V，第 41 页，命题 4 (ii)，我们有 \(u(\mathrm{F}^{\circ}) \subset \mathrm{F}^{\circ}\) 与 \(u^{*}(\mathrm{F}^{\circ}) \subset \mathrm{F}^{\circ}\) 。由此得出 \(\widetilde{u}\) 的伴随是 \(\mathrm{F}^{\circ}\) 的由 \(u^{*}\) 诱导的自同态。由于 \(u\) 是正规的，自同态 \(\widetilde{u}\) 是正规的。

